% sections2/sec4_linear_base_ext.tex
\section{线性映射的张量积与基变换}
\label{sec:linear-base-ext}

固定环同态 $R \xrightarrow{f} S$.每个 $S$-模均可通过标量限制成为 $R$-模，而每个 $R$-模 $M$ 均具有基变换（标量扩张）$S \otimes_R M$（它是一个 $S$-模）.在第一篇中我们看到，$S \otimes_R M$ 在所有 $S$-模中满足普遍映射性质：从 $M$ 出发到达任意 $S$-模的 $R$-线性映射均可唯一延拓为从 $S \otimes_R M$ 出发的 $S$-线性映射.

在本节中，我们探讨基变换更为重要的另一面：它将两个 $R$-模之间的 $R$-线性映射 $M \xrightarrow{\varphi} M'$ 自然转化为 $S$-模之间的 $S$-线性映射.将映射 $\varphi$ 与 $S$ 作张量积，得到 $R$-线性映射
\[
  S \otimes_R M \xrightarrow{1 \otimes \varphi} S \otimes_R M'.
\]
容易验证该映射事实上是 $S$-线性的：$(1 \otimes \varphi)(st) = s(1 \otimes \varphi)(t)$.为了突显其对 $S$ 的依赖性，我们通常将该基变换后的线性映射记作 $\varphi_S$：
\begin{equation}
  \varphi_S \colon S \otimes_R M \to S \otimes_R M', \quad \varphi_S(s \otimes m) = s \otimes \varphi(m).
\end{equation}
由于 $(\id_M)_S = \id_{S \otimes_R M}$ 且 $(\varphi \circ \varphi')_S = \varphi_S \circ \varphi'_S$，基变换过程具有完整的函子性.

若 $R$-线性映射 $\varphi$ 是同构或满射，则其基变换 $\varphi_S$ 亦保持同构或满射.但若 $\varphi$ 是单射，$\varphi_S$ 却未必是单射（见第一篇及第二节的反例）.

\begin{theorem}\label{thm:invariance-rank-domain}
设 $R$ 为非零交换环.若 $R^m \cong R^n$ 为 $R$-模同构，则 $m = n$；若存在满线性映射 $R^m \twoheadrightarrow R^n$，则 $m \geq n$.
\end{theorem}

\begin{proof}
选取 $R$ 的极大理想 $\mathfrak{m}$.将映射与域 $R/\mathfrak{m}$ 作张量积，即可转化为域上向量空间的线性同构或满射，由向量空间维数理论即得 $m = n$ 或 $m \geq n$.
\end{proof}

\begin{theorem}\label{thm:pid-free-rank}
设 $R$ 为 PID，$M$ 为有限生成 $R$-模.将 $M$ 分解为不变量形式：
\[
  M \cong R^d \oplus R/(a_1) \oplus \dots \oplus R/(a_k) \quad (a_1 \mid \dots \mid a_k),
\]
则自由秩 $d$ 恰好等于 $\dim_K(K \otimes_R M)$（其中 $K$ 为 $R$ 的分式域）.因此自由秩 $d$ 由 $M$ 唯一定义.
\end{theorem}

\begin{theorem}[基变换保持矩阵表示]\label{thm:base-ext-matrix}
设 $M$ 与 $M'$ 为非零有限自由 $R$-模，$\varphi \colon M \to M'$ 为 $R$-线性映射.任取基 $\{e_j\}$ 与 $\{e'_i\}$，则 $S$-线性映射 $\varphi_S \colon S \otimes_R M \to S \otimes_R M'$ 关于基 $\{1 \otimes e_j\}$ 与 $\{1 \otimes e'_i\}$ 的矩阵表示，与 $\varphi$ 关于基 $\{e_j\}$ 与 $\{e'_i\}$ 的矩阵表示完全相同！
\end{theorem}

\begin{proof}
若 $\varphi(e_j) = \sum_i a_{ij} e'_i$，则
\[
  \varphi_S(1 \otimes e_j) = 1 \otimes \varphi(e_j) = \sum_i a_{ij} (1 \otimes e'_i),
\]
矩阵系数全然未变.
\end{proof}

\begin{theorem}\label{thm:s-tensor-hom}
设 $M$ 与 $M'$ 为 $R$-模.存在唯一的 $S$-线性映射
\[
  S \otimes_R \Hom_R(M, M') \to \Hom_S(S \otimes_R M, S \otimes_R M'), \quad s \otimes \varphi \mapsto s\varphi_S.
\]
当 $M$ 与 $M'$ 为有限自由 $R$-模时，该映射是 $S$-模同构.特别地，取 $M' = R$，有自然同构
\[
  S \otimes_R (M^{\vee_R}) \cong (S \otimes_R M)^{\vee_S}.
\]
\end{theorem}

\begin{theorem}[整环上自由子模单射张量幂保持定理]\label{thm:domain-free-tensor-power-injective}
设 $R$ 为整环，$\varphi \colon M \hookrightarrow N$ 为单射线性映射.若 $M$ 是自由模，则对任意 $k \geq 1$，张量幂映射
\[
  \varphi^{\otimes k} \colon M^{\otimes k} \to N^{\otimes k}
\]
必为单射.
\end{theorem}

\begin{proof}
考虑交换图：
\begin{center}
\begin{tikzcd}[row sep=large, column sep=large]
M^{\otimes k} \arrow[r, "\varphi^{\otimes k}"] \arrow[d, hook] & N^{\otimes k} \arrow[d] \\
K \otimes_R M^{\otimes k} \arrow[r, "1 \otimes \varphi^{\otimes k}"'] \arrow[d, "\cong"'] & K \otimes_R N^{\otimes k} \arrow[d, "\cong"] \\
(K \otimes_R M)^{\otimes_K k} \arrow[r, "\varphi_K^{\otimes k}"'] & (K \otimes_R N)^{\otimes_K k}
\end{tikzcd}
\end{center}
由分式域 $K$ 的平坦性，$\varphi_K \colon K \otimes_R M \to K \otimes_R N$ 是向量空间的单射，故其张量幂 $\varphi_K^{\otimes k}$ 必单.又因 $M^{\otimes k}$ 是自由模（无扭），左侧坚线复合是单射.由交换图追逐即知顶层 $\varphi^{\otimes k}$ 必为单射.
\end{proof}
